% Seção 5 de replicacao_secao4.tex: as duas regressões do MQ2E, passo a passo.
% Números de gerado/numeros.tex (R/05_documento.R).
\section{As duas regressões, passo a passo}\label{sec:duas}

\subsection{Por que MQO não serve}
Em~\eqref{eq:estrutural}, $X^h_t$ contém a reação do Copom ao que ele vê da economia, e o erro
$u^h_t$ contém o que move o crédito além da Selic. Para ver o viés, residualize tudo nos controles
$W$ (Frisch--Waugh--Lovell): seja $M_W = I - W(W'W)^{-1}W'$ e $\tilde a = M_W a$ o resíduo de $a$
regredido em $W$. O MQO do coeficiente da Selic é
\begin{equation}\label{eq:mqo}
\hat\beta^{\,MQO} = \frac{\tilde X'\tilde Y}{\tilde X'\tilde X}
                  = \beta + \frac{\tilde X'u}{\tilde X'\tilde X}.
\end{equation}
Se o Copom reage ao que está em $u$, então $\tilde X'u \neq 0$ mesmo em amostras grandes, e o MQO
erra por esse termo. A variável instrumental troca, no numerador e no denominador, o $\tilde X$ por
uma variável que não reage a $u$.

\subsection{As condições da variável instrumental}
Um instrumento $z_t$ precisa de duas condições, dados os controles $W_t$:
\begin{enumerate}
  \item \textbf{Relevância}: $z_t$ move $X^h_t$, isto é, $\tilde z'\tilde X \neq 0$. Testável, pelo primeiro estágio.
  \item \textbf{Exogeneidade e exclusão}: $\operatorname{E}[z_t\,u^h_t\mid W_t]=0$; $z_t$ só afeta
        $Y^h_t$ por meio de $X^h_t$. \textbf{Não testável} com um instrumento e uma variável endógena:
        ``the exclusion or exogeneity condition cannot be tested'' e ``the reliability of an instrument
        relies on argumentation, not on empirical testing'' (Verbeek, 2017, p.~154).
\end{enumerate}
O BC escreve as duas como hipótese (p.~45--46) e as defende numa frase: as surpresas ``são
plausivelmente ortogonais às condições do mercado de crédito'' (p.~45). A ameaça principal é o efeito
informação (p.~42--43; Nakamura e Steinsson, 2018, p.~76).

\subsection{As três regressões}
O MQ2E costuma ser apresentado como duas regressões; vale escrever também a terceira, que liga as
duas. Os controles $W_t$ entram nas três.
\begin{align}
\text{Primeiro estágio:}\quad & X^h_t = \pi_h\, z_t + W_t'\delta_h + v^h_t, \label{eq:fs}\\
\text{Forma reduzida:}\quad & Y^h_t = \rho_h\, z_t + W_t'\lambda_h + e^h_t, \label{eq:rf}\\
\text{Segundo estágio:}\quad & Y^h_t = \beta_h\, \widehat{X}^h_t + W_t'\gamma_h + \text{erro},
\quad \widehat{X}^h_t = \hat\pi_h z_t + W_t'\hat\delta_h. \label{eq:ss}
\end{align}
O primeiro estágio mede quanto a surpresa move a Selic; a forma reduzida, quanto a surpresa move o
crédito; o segundo estágio usa só a parte da Selic que a surpresa (e os controles) explicam.

\subsection{O estimador em uma linha}
Com um instrumento e uma variável endógena, depois de residualizar nos controles:
\begin{align}
\text{1º estágio:}\quad & \hat\pi = \frac{\tilde z'\tilde X}{\tilde z'\tilde z}, \label{eq:pi}\\
\text{forma reduzida:}\quad & \hat\rho = \frac{\tilde z'\tilde Y}{\tilde z'\tilde z}, \label{eq:rho}\\
\text{MQ2E:}\quad & \hat\beta = \frac{\tilde z'\tilde Y}{\tilde z'\tilde X} = \frac{\hat\rho}{\hat\pi}. \label{eq:b2sls}
\end{align}
Compare com~\eqref{eq:mqo}: o MQ2E é o MQO com $\tilde z$ no lugar de $\tilde X$ nos dois
produtos. Substituindo $\tilde Y = \beta\tilde X + \tilde u$,
\begin{equation}\label{eq:erro_vi}
\hat\beta - \beta = \frac{\tilde z'u}{\tilde z'\tilde X},
\end{equation}
que vai a zero se $z$ não anda junto com $u$ (exogeneidade) e se o denominador não é perto de zero
(relevância). Se $\tilde z'\tilde X\approx0$, qualquer correlação pequena entre $z$ e $u$ explode:
é o problema do instrumento fraco.

\textbf{Em matrizes.} Com $\bar X=[X\ \ W]$, $\bar Z=[z\ \ W]$ e
$P_{\bar Z}=\bar Z(\bar Z'\bar Z)^{-1}\bar Z'$, o MQ2E é
\[
\hat\theta=(\bar X'P_{\bar Z}\bar X)^{-1}\bar X'P_{\bar Z}\,Y,
\qquad\text{e, no caso exatamente identificado,}\qquad
\hat\theta=(\bar Z'\bar X)^{-1}\bar Z'Y.
\]
A coordenada de $X$ nesse vetor é exatamente~\eqref{eq:b2sls}. Os ``dois estágios'' são a leitura de $P_{\bar Z}\bar X$: primeiro
projeta-se $X$ no espaço de $(z,W)$, depois regride-se $Y$ na projeção.

%-------------------------------------------------------------------------------------
\subsection{Exemplo 1 com todos os números: taxa do crédito livre, $h=18$}

\paragraph{Passo 1: montar as variáveis.} $Y^{18}_t$ e $X^{18}_t$ como em~\eqref{eq:Y}
e~\eqref{eq:X}; $z_t$ é a surpresa do mês; $W_t$ tem \KwTL{} colunas: constante, $\Delta y_{t-1}$,
$\Delta y_{t-2}$, $\Delta i_{t-1}$, $\Delta i_{t-2}$. São $N=\NTL$ meses, de \IniTL{} a 12/2023.

\paragraph{Passo 2: residualizar nos controles.} Regride-se $z$, $X$ e $Y$, cada um, em $W$ e
guardam-se os resíduos. Alguns meses:
\begin{table}[H]\centering\small
\caption{Variáveis residualizadas nos controles (taxa do crédito livre)}
\begin{tabular}{lrrrr}
\toprule
Mês $t$ & $\tilde z_t$ & $\tilde X_t$ & $\tilde Y_t$ & $g_t=\tilde z_t\hat u_t$ \\
\midrule
01/2021 & 0{,}010 & 4{,}907 & 6{,}927 & \ensuremath{-}0{,}002 \\
02/2021 & \ensuremath{-}0{,}003 & 5{,}848 & 4{,}831 & 0{,}012 \\
03/2021 & 0{,}242 & 7{,}507 & 6{,}803 & \ensuremath{-}1{,}001 \\
04/2021 & \ensuremath{-}0{,}003 & 5{,}240 & 3{,}808 & 0{,}010 \\
05/2021 & \ensuremath{-}0{,}005 & 4{,}927 & 2{,}328 & 0{,}026 \\
06/2021 & \ensuremath{-}0{,}006 & 3{,}816 & 1{,}890 & 0{,}022 \\
\bottomrule
\end{tabular}

\end{table}
Somando os produtos em todos os \NTL{} meses:
\[
\tilde z'\tilde z = \ZZTL,\qquad \tilde z'\tilde X = \ZXTL,\qquad \tilde z'\tilde Y = \ZYTL.
\]

\paragraph{Passo 3: primeiro estágio.} Por~\eqref{eq:pi}, $\hat\pi=\ZXTL/\ZZTL=\PmanTL$: uma surpresa
de 1~p.p. vem acompanhada de uma Selic média, nos 19 meses seguintes, $\PmanTL$~p.p. acima da do
mês anterior, já descontados os controles. As surpresas da amostra são de 0,25 ou 0,50~p.p., então
uma surpresa típica corresponde a um quarto disso; o coeficiente é grande porque as surpresas
costumam cair no começo de ciclos de juros. É o coeficiente de $z_t$ na regressão completa:
\begin{table}[H]\centering\small
\caption{Primeiro estágio, equação~\eqref{eq:fs}: $X^{18}_t$ na surpresa e nos controles}
\begin{tabular}{lrrr}
\toprule
Regressor & Coeficiente & Erro-padrão NW & $t$ \\
\midrule
$z_t$ (surpresa do Copom) & 6{,}514 & 1{,}645 & 3{,}96 \\
$\Delta y_{t-1}$ & \ensuremath{-}0{,}164 & 0{,}152 & \ensuremath{-}1{,}07 \\
$\Delta y_{t-2}$ & \ensuremath{-}0{,}434 & 0{,}185 & \ensuremath{-}2{,}34 \\
$\Delta i_{t-1}$ & 3{,}932 & 0{,}734 & 5{,}36 \\
$\Delta i_{t-2}$ & 1{,}202 & 0{,}491 & 2{,}45 \\
Constante & 0{,}044 & 0{,}528 & 0{,}08 \\
\midrule
\multicolumn{4}{l}{$N$ = 151 meses (06/2011 a 12/2023); $R^2$ = 0{,}395} \\
\bottomrule
\end{tabular}

\end{table}
O F efetivo é o quadrado da estatística $t$ de $z_t$ com erro de Newey-West: $\FefTL$. Acima de 10,
a regra prática (Verbeek, 2017, p.~153).

\paragraph{Passo 4: forma reduzida.} Por~\eqref{eq:rho}, $\hat\rho=\ZYTL/\ZZTL=\RmanTL$: a mesma
surpresa vem acompanhada de uma taxa do crédito livre $\RmanTL$~p.p. mais alta, no mesmo sentido.
\begin{table}[H]\centering\small
\caption{Forma reduzida, equação~\eqref{eq:rf}: $Y^{18}_t$ na surpresa e nos controles}
\begin{tabular}{lrrr}
\toprule
Regressor & Coeficiente & Erro-padrão NW & $t$ \\
\midrule
$z_t$ (surpresa do Copom) & 9{,}498 & 3{,}394 & 2{,}80 \\
$\Delta y_{t-1}$ & \ensuremath{-}0{,}130 & 0{,}291 & \ensuremath{-}0{,}45 \\
$\Delta y_{t-2}$ & \ensuremath{-}0{,}423 & 0{,}275 & \ensuremath{-}1{,}54 \\
$\Delta i_{t-1}$ & 6{,}753 & 1{,}018 & 6{,}64 \\
$\Delta i_{t-2}$ & 3{,}437 & 0{,}825 & 4{,}17 \\
Constante & 0{,}272 & 0{,}847 & 0{,}32 \\
\midrule
\multicolumn{4}{l}{$N$ = 151 meses (06/2011 a 12/2023); $R^2$ = 0{,}493} \\
\bottomrule
\end{tabular}

\end{table}

\paragraph{Passo 5: segundo estágio.} Por~\eqref{eq:b2sls}, $\hat\beta=\ZYTL/\ZXTL=\BmanTL$, igual
a $\hat\rho/\hat\pi=\RmanTL/\PmanTL$ e igual ao coeficiente de $\widehat X$ na regressão do
segundo estágio:
\begin{table}[H]\centering\small
\caption{Segundo estágio, equação~\eqref{eq:ss}: $Y^{18}_t$ no valor previsto e nos controles}
\begin{tabular}{lrrr}
\toprule
Regressor & Coeficiente & Erro-padrão NW & $t$ \\
\midrule
$\widehat{X}^h_t$ (valor previsto no 1º estágio) & 1{,}458 & 0{,}521 & 2{,}80 \\
$\Delta y_{t-1}$ & 0{,}108 & 0{,}316 & 0{,}34 \\
$\Delta y_{t-2}$ & 0{,}209 & 0{,}282 & 0{,}74 \\
$\Delta i_{t-1}$ & 1{,}020 & 2{,}397 & 0{,}43 \\
$\Delta i_{t-2}$ & 1{,}684 & 1{,}214 & 1{,}39 \\
Constante & 0{,}207 & 0{,}844 & 0{,}25 \\
\midrule
\multicolumn{4}{l}{$N$ = 151 meses (06/2011 a 12/2023); $R^2$ = 0{,}493} \\
\bottomrule
\end{tabular}

\end{table}
Lê-se: uma Selic 1~p.p. mais alta, em média, nos 19 meses deixa a taxa média do crédito livre
$\BetaTL$~p.p. mais alta no mesmo período. O BC obtém 1,05 (p.~47). O $R^2$ do segundo estágio é
igual ao da forma reduzida porque $\widehat X$ é combinação de $z$ e $W$: as duas regressões usam o
mesmo espaço de regressores.

\paragraph{Passo 6: que variação da Selic a regressão usou.} A tabela a seguir põe lado a lado
$X_t$ (o que a Selic fez) e $\widehat X_t$ (a parte que a surpresa e os controles explicam):
\begin{table}[H]\centering\small
\caption{O que entra em cada regressão, em volta do início do ciclo de alta de 2021}
\resizebox{\textwidth}{!}{\begin{tabular}{lrrrrrrr}
\toprule
Mês $t$ & $z_t$ & $X^{18}_t$ & $\widehat{X}^{18}_t$ & $Y^{18}_t$ & $\hat\beta X_t + W_t'\hat\gamma$ & $\hat u_t$ (certo) & $\hat e_t$ (2º estágio à mão) \\
\midrule
01/2021 & 0{,}00 & 5{,}15 & 0{,}31 & 7{,}34 & 7{,}57 & \ensuremath{-}0{,}23 & 6{,}83 \\
02/2021 & 0{,}00 & 5{,}77 & \ensuremath{-}0{,}10 & 5{,}08 & 8{,}77 & \ensuremath{-}3{,}69 & 4{,}86 \\
03/2021 & 0{,}25 & 6{,}39 & 0{,}45 & 5{,}96 & 10{,}10 & \ensuremath{-}4{,}14 & 4{,}51 \\
04/2021 & 0{,}00 & 6{,}66 & 1{,}40 & 6{,}40 & 10{,}23 & \ensuremath{-}3{,}83 & 3{,}83 \\
05/2021 & 0{,}00 & 6{,}82 & 1{,}86 & 6{,}41 & 11{,}27 & \ensuremath{-}4{,}86 & 2{,}38 \\
06/2021 & 0{,}00 & 6{,}72 & 2{,}87 & 7{,}77 & 11{,}44 & \ensuremath{-}3{,}67 & 1{,}95 \\
\bottomrule
\end{tabular}
}
\end{table}
Em 2021 a Selic média dos 19 meses seguintes ficou cerca de 5 a 7~p.p. acima da do mês anterior
($X_t$), mas $\widehat X_t$ fica perto de zero até março: o mercado já esperava quase todo o ciclo.
Depois de março, $\widehat X_t$ sobe porque as altas já feitas entram nos controles
($\Delta i_{t-1}$, $\Delta i_{t-2}$), não por novas surpresas. A VI ignora a parte antecipada e
estima $\beta$ só com a variação ligada à surpresa. É isso que a protege da reação do
Copom à economia, e é também por isso que ela tem pouca informação (\Nmesessurp{} meses com
$z_t\neq0$ na amostra).

\paragraph{Passo 7: os resíduos certos.} O resíduo do modelo usa a Selic de verdade:
$\hat u_t = Y_t - \hat\beta X_t - W_t'\hat\gamma$. Quem roda o segundo estágio à mão obtém
$\hat e_t = Y_t - \hat\beta \widehat X_t - W_t'\hat\gamma$, com os mesmos coeficientes e o resíduo
errado (as duas últimas colunas da tabela acima). A soma dos quadrados dos resíduos certos é
$\SSRokTL$; a dos errados, $\SSRssTL$. O erro-padrão que sai do segundo estágio à mão ($\SEssTL$) é,
por isso, diferente do correto ($\SEivTL$).

\paragraph{Passo 8: o erro-padrão de Newey-West, à mão.} De~\eqref{eq:erro_vi},
$\hat\beta-\beta = \sum_t g_t / \tilde z'\tilde X$ com $g_t=\tilde z_t u_t$ (última coluna da
tabela do passo 2, com $\hat u_t$ no lugar de $u_t$). Então
\begin{equation}\label{eq:nw}
\widehat{\operatorname{Var}}(\hat\beta) = \frac{\hat S}{(\tilde z'\tilde X)^2},\qquad
\hat S = \sum_t g_t^2 + 2\sum_{j=1}^{L}\Bigl(1-\frac{j}{L+1}\Bigr)\sum_t g_t\,g_{t-j}.
\end{equation}
O primeiro termo de $\hat S$ trata a heterocedasticidade; o segundo, a correlação entre meses
próximos, com pesos que caem linearmente (Bartlett). Usa-se $L=h+1=\LnwTL$ porque $Y^{18}_t$ e
$Y^{18}_{t+1}$ compartilham 18 dos 19 meses: os erros de meses vizinhos são correlacionados por
construção. No exemplo, $\hat S=\SmanTL$ e
\[
\widehat{\text{ep}}(\hat\beta)=\frac{\sqrt{\hat S}}{|\tilde z'\tilde X|}=\frac{\sqrt{\SmanTL}}{\ZXTL}=\SEmanTL,
\]
o mesmo valor do pacote (\texttt{sandwich::NeweyWest}). Para comparar: sem o termo de
autocorrelação, o erro-padrão seria $\SEhcTL$; supondo erros homocedásticos e independentes,
$\SEhomTL$. O intervalo de 95\% é $\hat\beta\pm1{,}96\,\widehat{\text{ep}} = [\IClTL;\ \IChTL]$.

%-------------------------------------------------------------------------------------
\subsection{Exemplo 2: concessões a empresas com recursos direcionados, $h=18$}
Mesmos passos; agora $W_t$ tem \KwCD{} colunas (as das taxas mais duas defasagens da variação da
desocupação, do log do IPCA e do log do IBC-Br, e a dummy de pandemia), e a amostra começa em
\IniCD{} porque a desocupação da PNAD Contínua só existe desde 2012 ($N=\NCD$).

\begin{table}[H]\centering\small
\caption{Primeiro estágio: $X^{18}_t$ na surpresa (amostra das concessões)}
\begin{tabular}{lrrr}
\toprule
Regressor & Coeficiente & Erro-padrão NW & $t$ \\
\midrule
$z_t$ (surpresa do Copom) & 6{,}568 & 1{,}269 & 5{,}18 \\
Dummy de pandemia & 2{,}823 & 1{,}124 & 2{,}51 \\
$\Delta y_{t-1}$ & 0{,}009 & 0{,}007 & 1{,}23 \\
$\Delta y_{t-2}$ & 0{,}006 & 0{,}006 & 0{,}98 \\
$\Delta i_{t-1}$ & 1{,}827 & 0{,}472 & 3{,}87 \\
$\Delta i_{t-2}$ & \ensuremath{-}0{,}244 & 0{,}679 & \ensuremath{-}0{,}36 \\
$\Delta u_{t-1}$ & \ensuremath{-}0{,}221 & 0{,}575 & \ensuremath{-}0{,}38 \\
$\Delta u_{t-2}$ & \ensuremath{-}0{,}522 & 0{,}568 & \ensuremath{-}0{,}92 \\
$\Delta p_{t-1}$ & 1{,}267 & 0{,}377 & 3{,}36 \\
$\Delta p_{t-2}$ & 1{,}115 & 0{,}400 & 2{,}79 \\
$\Delta a_{t-1}$ & 0{,}029 & 0{,}081 & 0{,}36 \\
$\Delta a_{t-2}$ & \ensuremath{-}0{,}022 & 0{,}050 & \ensuremath{-}0{,}45 \\
Constante & \ensuremath{-}1{,}612 & 0{,}765 & \ensuremath{-}2{,}11 \\
\midrule
\multicolumn{4}{l}{$N$ = 139 meses (06/2012 a 12/2023); $R^2$ = 0{,}571} \\
\bottomrule
\end{tabular}

\end{table}
\begin{table}[H]\centering\small
\caption{Forma reduzida: $Y^{18}_t$ (log das concessões $\times100$) na surpresa}
\begin{tabular}{lrrr}
\toprule
Regressor & Coeficiente & Erro-padrão NW & $t$ \\
\midrule
$z_t$ (surpresa do Copom) & \ensuremath{-}17{,}714 & 26{,}312 & \ensuremath{-}0{,}67 \\
Dummy de pandemia & 9{,}912 & 9{,}499 & 1{,}04 \\
$\Delta y_{t-1}$ & \ensuremath{-}0{,}568 & 0{,}098 & \ensuremath{-}5{,}82 \\
$\Delta y_{t-2}$ & \ensuremath{-}0{,}289 & 0{,}122 & \ensuremath{-}2{,}36 \\
$\Delta i_{t-1}$ & 0{,}604 & 8{,}703 & 0{,}07 \\
$\Delta i_{t-2}$ & \ensuremath{-}1{,}647 & 8{,}492 & \ensuremath{-}0{,}19 \\
$\Delta u_{t-1}$ & \ensuremath{-}5{,}893 & 9{,}787 & \ensuremath{-}0{,}60 \\
$\Delta u_{t-2}$ & \ensuremath{-}21{,}397 & 12{,}576 & \ensuremath{-}1{,}70 \\
$\Delta p_{t-1}$ & \ensuremath{-}14{,}426 & 6{,}488 & \ensuremath{-}2{,}22 \\
$\Delta p_{t-2}$ & 0{,}253 & 5{,}827 & 0{,}04 \\
$\Delta a_{t-1}$ & \ensuremath{-}2{,}657 & 1{,}940 & \ensuremath{-}1{,}37 \\
$\Delta a_{t-2}$ & \ensuremath{-}2{,}334 & 1{,}317 & \ensuremath{-}1{,}77 \\
Constante & \ensuremath{-}0{,}033 & 7{,}836 & \ensuremath{-}0{,}00 \\
\midrule
\multicolumn{4}{l}{$N$ = 139 meses (06/2012 a 12/2023); $R^2$ = 0{,}284} \\
\bottomrule
\end{tabular}

\end{table}
\begin{table}[H]\centering\small
\caption{Segundo estágio: $Y^{18}_t$ no valor previsto}
\begin{tabular}{lrrr}
\toprule
Regressor & Coeficiente & Erro-padrão NW & $t$ \\
\midrule
$\widehat{X}^h_t$ (valor previsto no 1º estágio) & \ensuremath{-}2{,}697 & 4{,}006 & \ensuremath{-}0{,}67 \\
Dummy de pandemia & 17{,}527 & 13{,}937 & 1{,}26 \\
$\Delta y_{t-1}$ & \ensuremath{-}0{,}545 & 0{,}091 & \ensuremath{-}6{,}00 \\
$\Delta y_{t-2}$ & \ensuremath{-}0{,}272 & 0{,}115 & \ensuremath{-}2{,}36 \\
$\Delta i_{t-1}$ & 5{,}531 & 11{,}082 & 0{,}50 \\
$\Delta i_{t-2}$ & \ensuremath{-}2{,}304 & 8{,}325 & \ensuremath{-}0{,}28 \\
$\Delta u_{t-1}$ & \ensuremath{-}6{,}490 & 10{,}011 & \ensuremath{-}0{,}65 \\
$\Delta u_{t-2}$ & \ensuremath{-}22{,}805 & 12{,}206 & \ensuremath{-}1{,}87 \\
$\Delta p_{t-1}$ & \ensuremath{-}11{,}009 & 9{,}437 & \ensuremath{-}1{,}17 \\
$\Delta p_{t-2}$ & 3{,}261 & 6{,}937 & 0{,}47 \\
$\Delta a_{t-1}$ & \ensuremath{-}2{,}578 & 1{,}956 & \ensuremath{-}1{,}32 \\
$\Delta a_{t-2}$ & \ensuremath{-}2{,}394 & 1{,}318 & \ensuremath{-}1{,}82 \\
Constante & \ensuremath{-}4{,}382 & 9{,}068 & \ensuremath{-}0{,}48 \\
\midrule
\multicolumn{4}{l}{$N$ = 139 meses (06/2012 a 12/2023); $R^2$ = 0{,}284} \\
\bottomrule
\end{tabular}

\end{table}

\textbf{As contas.} $\tilde z'\tilde z=\ZZCD$, $\tilde z'\tilde X=\ZXCD$, $\tilde z'\tilde Y=\ZYCD$.
Primeiro estágio $\hat\pi=\PmanCD$; forma reduzida $\hat\rho=\RmanCD$; $\hat\beta=\ZYCD/\ZXCD=\BmanCD$:
por ponto de Selic, as concessões direcionadas a empresas ficam $\BetaCD$\% menores em média nos 18
meses. O BC reporta $-5{,}3$\%, significante (p.~51). Erro-padrão de Newey-West $\SEmanCD$ (sem
autocorrelação, $\SEhcCD$; homocedástico, $\SEhomCD$); intervalo $[\IClCD;\ \IChCD]$, com o zero
dentro. F efetivo $\FefCD$.

Note que a forma reduzida já conta a história toda: a surpresa tem coeficiente $\PiRFCD$ com
estatística $t$ pequena, isto é, não há relação detectável entre surpresas e concessões
direcionadas a empresas nesta amostra. Como $\hat\beta=\hat\rho/\hat\pi$ e o primeiro estágio é
forte, o MQ2E só reescala essa forma reduzida imprecisa em unidades de Selic.

\subsection{MQO contra VI}
Sem instrumento, o MQO de~\eqref{eq:estrutural} dá $\BetaTLols$ para a taxa do crédito livre
(contra $\BetaTL$ por VI) e $\BetaCDols$ para as concessões direcionadas a empresas (contra
$\BetaCD$). Por~\eqref{eq:mqo}, a diferença é uma estimativa do viés $\tilde X'u/\tilde X'\tilde X$,
se o instrumento for válido.
